\section*{Capítulo \ref{ch:g11:functional-groups} --- Grupos funcionales y familias}

\begin{solution}{exo:g11:functional-groups:1}
Hidroxilo, \omterm{def:g11:functional-groups:alcohol}{alcohol}; carbonilo en el interior de la cadena, \omterm{def:g11:functional-groups:carbonyl}{cetona};
carboxilo, \omterm{def:g11:functional-groups:carboxylic-acid}{ácido carboxílico}; grupo \omterm{def:g11:functional-groups:carboxylic-acid}{éster}, \omterm{def:g11:functional-groups:carboxylic-acid}{éster}; grupo amino
\ce{-NH2}, \omterm{def:g11:functional-groups:amine}{amina}.
\end{solution}

\begin{solution}{exo:g11:functional-groups:2}
Metanol; propan-1-ol; ácido metanoico; propanal; propanona.
\end{solution}

\begin{solution}{exo:g11:functional-groups:3}
\omterm{def:g11:functional-groups:alcohol}{Alcohol}; \omterm{def:g11:functional-groups:carbonyl}{aldehído}; \omterm{def:g11:functional-groups:carboxylic-acid}{éster}; \omterm{def:g11:functional-groups:amine}{amida}; \omterm{def:g11:functional-groups:halogenoalkane}{halogenoalcano}; \omterm{def:g11:functional-groups:halogenoalkane}{alqueno}.
\end{solution}

\begin{solution}{exo:g11:functional-groups:4}
Propanal \ce{CH3-CH2-CHO}, propanona \ce{CH3-CO-CH3}. El propanal es el
\omterm{def:g11:functional-groups:carbonyl}{aldehído}: el carbono de su \ce{C=O} está en el extremo de la cadena,
unido a un \omterm{def:g7:atoms-and-molecules:atom}{átomo} de hidrógeno; en la propanona está entre dos \omterm{def:g7:atoms-and-molecules:atom}{átomos} de
carbono.
\end{solution}

\begin{solution}{exo:g11:functional-groups:5}
Primario (el carbono que lleva el \ce{-OH} tiene un carbono vecino);
secundario (dos); terciario (tres).
\end{solution}

\begin{solution}{exo:g11:functional-groups:6}
\setchemfig{atom sep=1.5em}
Pentan-3-ol \chemfig{-[:30]-[:-30](-[:-90]OH)-[:30]-[:-30]};
ácido 2-metilbutanoico \chemfig{-[:30]-[:-30](-[:-90])-[:30](=[:90]O)-[:-30]OH};
hexan-2-ona \chemfig{-[:30](=[:90]O)-[:-30]-[:30]-[:-30]-[:30]};
etanoato de propilo \chemfig{-[:30](=[:90]O)-[:-30]O-[:30]-[:-30]-[:30]}.
\end{solution}

\begin{solution}{exo:g11:functional-groups:7}
\ce{C3H6O}. Sí, son \omterm{def:g11:organic-skeletons:isomer}{isómeros}, pero con grupos distintos: un \omterm{def:g11:functional-groups:carbonyl}{aldehído} y
una \omterm{def:g11:functional-groups:carbonyl}{cetona} (los dos grupos carbonilo, en sitios distintos).
\end{solution}

\begin{solution}{exo:g11:functional-groups:8}
\ce{HCOO-CH2-CH3}, metanoato de etilo; \ce{CH3-COO-CH2-CH2-CH3},
etanoato de propilo.
\end{solution}

\begin{solution}{exo:g11:functional-groups:9}
Etanoato de metilo (ácido etanoico, metanol); propanoato de etilo
(ácido propanoico, etanol); metanoato de metilo (ácido metanoico,
metanol).
\end{solution}

\begin{solution}{exo:g11:functional-groups:10}
\setchemfig{atom sep=1.5em}
\begin{center}
\small
\begin{tabular}{cccc}
\chemfig{-[:30]-[:-30]-[:30]-[:-30]OH} &
\chemfig{-[:30](-[:90]OH)-[:-30]-[:30]} &
\chemfig{-[:30](-[:90])-[:-30]-[:30]OH} &
\chemfig{-[:0](-[:90]OH)(-[:-90])-[:0]} \\
butan-1-ol & butan-2-ol & 2-metilpropan-1-ol & 2-metilpropan-2-ol \\
primario & secundario & primario & terciario
\end{tabular}
\end{center}
\end{solution}

\begin{solution}{exo:g11:functional-groups:11}
Una \omterm{def:g11:functional-groups:amine}{amida} (etanamida). \emph{-eno}. Los \omterm{def:g11:functional-groups:carboxylic-acid}{ácidos carboxílicos} (\ce{C=O} y
\ce{-OH} sobre el mismo carbono), los \omterm{def:g11:functional-groups:carboxylic-acid}{ésteres} (\ce{C=O} y \ce{-O-}) y
las \omterm{def:g11:functional-groups:amine}{amidas} (\ce{C=O} y \ce{N}): tres familias.
\end{solution}

\begin{solution}{exo:g11:functional-groups:12}
Aspirina: un grupo carboxilo (\omterm{def:g11:functional-groups:carboxylic-acid}{ácido carboxílico}) y un grupo \omterm{def:g11:functional-groups:carboxylic-acid}{éster}.
Paracetamol: un \ce{-OH} fenol sobre el anillo y un grupo \omterm{def:g11:functional-groups:amine}{amida}
\ce{NH-CO}.
\end{solution}

\begin{solution}{exo:g11:functional-groups:13}
\ce{C2H6O}. Solo el etanol, con su \ce{O-H}, forma \omterm{def:g11:polarity-and-cohesion:hydrogen-bond}{enlaces de hidrógeno}
entre sus \omterm{def:g7:atoms-and-molecules:molecule}{moléculas} (el metoximetano no tiene hidrógeno sobre su
oxígeno). El etanol hierve a mayor temperatura, a \qty{78}{\celsius}; el
metoximetano, a \qty{-25}{\celsius}.
\end{solution}

\begin{solution}{exo:g11:functional-groups:14}
Un grupo hidroxilo y un grupo carboxilo. La cadena tiene tres carbonos;
el carbono del ácido es el carbono 1 y el \ce{-OH} está en el carbono 2:
ácido 2-hidroxipropanoico. La glicina: ácido 2-aminoetanoico (que a
menudo se escribe ácido aminoetanoico).
\end{solution}

\begin{solution}{exo:g11:functional-groups:15}
Etanol \ce{CH3-CH2-OH}, \ce{C2H6O}, \omterm{def:g11:functional-groups:alcohol}{alcohol}; etanal \ce{CH3-CHO},
\ce{C2H4O}, \omterm{def:g11:functional-groups:carbonyl}{aldehído}; ácido etanoico \ce{CH3-COOH}, \ce{C2H4O2}, \omterm{def:g11:functional-groups:carboxylic-acid}{ácido
carboxílico}. Del etanol al etanal se pierden dos \omterm{def:g7:atoms-and-molecules:atom}{átomos} de hidrógeno;
del etanal al ácido etanoico se gana un \omterm{def:g7:atoms-and-molecules:atom}{átomo} de oxígeno.
\end{solution}

\begin{solution}{pb:g11:functional-groups:1}
\textbf{1.} A: grupo \omterm{def:g11:functional-groups:carboxylic-acid}{éster}; B: hidroxilo; C: carboxilo; D: carbonilo en
el extremo de la cadena; E: grupo \omterm{def:g11:functional-groups:carboxylic-acid}{éster}.

\textbf{2.} A \omterm{def:g11:functional-groups:carboxylic-acid}{éster}, B \omterm{def:g11:functional-groups:alcohol}{alcohol}, C \omterm{def:g11:functional-groups:carboxylic-acid}{ácido carboxílico}, D \omterm{def:g11:functional-groups:carbonyl}{aldehído}, E
\omterm{def:g11:functional-groups:carboxylic-acid}{éster}.

\textbf{3.} A (plátano) y E (piña): olores afrutados.

\textbf{4.} Primario: el carbono que lleva el \ce{-OH} está unido a un
carbono.

\textbf{5.} \ce{C6H12O}; por ejemplo, la hexan-2-ona,
\ce{CH3-CO-CH2-CH2-CH2-CH3}.

\textbf{6.} La cadena que contiene el \ce{C-OH} tiene cuatro carbonos,
numerados desde el \ce{-OH}, con un metilo en el carbono 3:
3-metilbutan-1-ol.

\textbf{7.} C: ácido etanoico; D: hexanal.

\textbf{8.} Butanoato de etilo, relacionado con el ácido butanoico y el
etanol.

\textbf{9.} \emph{Etanoato}: la parte ácida \ce{CH3-CO}, del ácido
etanoico; \emph{3-metilbutilo}: la parte alcohólica, del
3-metilbutan-1-ol.

\textbf{10.} C (ácido etanoico) y B (3-metilbutan-1-ol).

\textbf{11.} A \ce{C7H14O2}; B \ce{C5H12O}; C \ce{C2H4O2}.

\textbf{12.} $\ce{C2H4O2} + \ce{C5H12O}$ contiene \ce{C7H16O3}, un
\ce{H2O} más que \ce{C7H14O2}: agua.

\textbf{13.} \ce{C2H4O2 + C5H12O -> C7H14O2 + H2O}: C 7 = 7, H 16 = 16,
O 3 = 3.

\textbf{14.} $M(\text{B}) = 5 \times 12.0 + 12 \times 1.0 + 16.0 =
\qty{88.0}{g/mol}$; $M(\text{C}) = \qty{60.0}{g/mol}$.

\textbf{15.} $60.0 + 88.0 - 18.0 = \qty{130.0}{g/mol}$.

\textbf{16.} $7 \times 12.0 + 14 \times 1.0 + 2 \times 16.0 =
\qty{130.0}{g/mol}$: la misma.

\textbf{17.} E, \ce{C6H12O2}: $6 \times 12.0 + 12 \times 1.0 + 2 \times
16.0 = \qty{116.0}{g/mol}$.

\textbf{18.} \ce{C7H14O2}: la misma fórmula que A, así que es un
\omterm{def:g11:organic-skeletons:isomer}{isómero}, con la misma \omterm{def:g10:the-mole:molar-mass}{masa molar}. Una \omterm{def:g10:the-mole:molar-mass}{masa molar}, por sí sola, no
permite distinguir \omterm{def:g11:organic-skeletons:isomer}{isómeros}.

\textbf{19.} A y el ácido heptanoico tienen los mismos \omterm{def:g7:atoms-and-molecules:atom}{átomos} pero
grupos distintos, un \omterm{def:g11:functional-groups:carboxylic-acid}{éster} y un \omterm{def:g11:functional-groups:carboxylic-acid}{ácido carboxílico}: uno huele a plátano,
y el otro, a rancio. El grupo, y el sitio donde está, deciden.

\textbf{20.} \qty{130.0}{g/mol}.
\end{solution}
